%!TEX TS-program = xelatex

\documentclass[c, dvipsnames]{beamer}  % [t], [c], или [b] --- вертикальное 
%\documentclass[handout, dvipsnames, c]{beamer} % Раздаточный материал (на слайдах всё сразу)
%выравнивание на слайдах (верх, центр, низ)
%\documentclass[handout, dvipsnames]{beamer} % Раздаточный материал (на слайдах всё сразу)
%\documentclass[aspectratio=169, dvipsnames]{beamer} % Соотношение сторон
\setbeamertemplate{navigation symbols}{}%remove navigation symbols

%\usetheme{Berkeley} % Тема оформленияLLL
%\usetheme{Bergen}
%\usetheme{CambridgeUS}
\usetheme{Boadilla}

\usecolortheme{crane} % Цветовая схема

%\useoutertheme{infolines} % Навигация 
%\useoutertheme{tree}
%\useoutertheme{miniframes}
%\useoutertheme{shadow}
%\useoutertheme{sidebar}
%\useoutertheme{smoothbars}
%\useoutertheme{smoothtree}
%\useoutertheme{split}
%\useoutertheme{default}


%\useinnertheme{circles}
\useinnertheme{rectangles}
%\useinnertheme{rounded}
%\useinnertheme{inmargin}


%%% Работа с русским языком
\usepackage[english,russian]{babel}   %% загружает пакет многоязыковой вёрстки
\usepackage{fontspec}      %% подготавливает загрузку шрифтов Open Type, True Type и др.
\defaultfontfeatures{Ligatures={TeX},Renderer=Basic}  %% свойства шрифтов по умолчанию
\setmainfont[Ligatures={TeX,Historic}]{Arial} %% задаёт основной шрифт документа
\setsansfont{Arial}                    %% задаёт шрифт без засечек
\setmonofont{Arial}
\usepackage{indentfirst}
\frenchspacing



%% Beamer по-русски
\newtheorem{rtheorem}{Теорема}
\newtheorem{rproof}{Доказательство}
\newtheorem{rexample}{Пример}

%%% Дополнительная работа с математикой
\usepackage{amsmath,amsfonts,amssymb,amsthm,mathtools} % AMS
\usepackage{icomma} % "Умная" запятая: $0,2$ --- число, $0, 2$ --- перечисление

%% Номера формул
\mathtoolsset{showonlyrefs=true} % Показывать номера только у тех формул, на которые есть \eqref{} в тексте.
%\usepackage{leqno} % Нумерация формул слева

%% Свои команды
\DeclareMathOperator{\sgn}{\mathop{sgn}}

%% Перенос знаков в формулах (по Львовскому)
\newcommand*{\hm}[1]{#1\nobreak\discretionary{}
{\hbox{$\mathsurround=0pt #1$}}{}}

%%% Работа с картинками
\usepackage{graphicx}  % Для вставки рисунков
\graphicspath{{images/}{images2/}}  % папки с картинками
\setlength\fboxsep{3pt} % Отступ рамки \fbox{} от рисунка
\setlength\fboxrule{1pt} % Толщина линий рамки \fbox{}
\usepackage{wrapfig} % Обтекание рисунков текстом

%%% Работа с таблицами
\usepackage{array,tabularx,tabulary,booktabs} % Дополнительная работа с таблицами
\usepackage{longtable}  % Длинные таблицы
\usepackage{multirow} % Слияние строк в таблице

%%% Программирование
\usepackage{etoolbox} % логические операторы

%%% Другие пакеты
\usepackage{lastpage} % Узнать, сколько всего страниц в документе.
%\usepackage{soul} % Модификаторы начертания
\usepackage{csquotes} % Еще инструменты для ссылок
\usepackage{multicol} % Несколько колонок


\usepackage{hyperref}
\usepackage{xcolor}
\hypersetup{        % Гиперссылки
    unicode=true,           % русские буквы в раздела PDF
    pdftitle={Заголовок},   % Заголовок
    pdfauthor={Автор},      % Автор
    pdfsubject={Тема},      % Тема
    pdfcreator={Создатель}, % Создатель
    pdfproducer={Производитель}, % Производитель
    pdfkeywords={keyword1} {key2} {key3}, % Ключевые слова
    colorlinks=true,        % false: ссылки в рамках; true: цветные ссылки
    linkcolor=,          % внутренние ссылки
    citecolor=green,        % на библиографию
    filecolor=magenta,      % на файлы
    urlcolor=blue           % на URL
} 

\usepackage{dcolumn}

%fffff3
\definecolor{backgr}{RGB}{146,26,29}
\definecolor{backgr1}{RGB}{230,43,37}
\definecolor{ex1}{RGB}{231,142,36}
\definecolor{ex2}{RGB}{249,155,28}
\definecolor{ex3}{RGB}{242,103,36}

\definecolor{red}{RGB}{230,43,37}
%\setbeamercolor{normal text}{fg=black,bg=backgr}
\setbeamercolor{frametitle}{bg=backgr,fg=white}
%\setbeamercolor{footline}{bg=backgr,fg=white}
%\setbeamercolor{normal text}{bg=yellow}
%\setbeamercolor{section in toc}{fg=yellow}
%\setbeamercolor{subsection in toc}{fg=blue}

% How to change colour of Navigation Bar in Beamer -  много интересного

%Пример команд, задающих внешний вид блока
\setbeamercolor{block title}{fg=white,bg=ex1}
\setbeamerfont{block title}{family=\sffamily}
\setbeamercolor{block body}{bg=white}
\setbeamertemplate{blocks}[rounded][shadow=fasle]
\setbeamercolor{title}{bg=backgr, fg=white}
\setbeamercolor{alerted text}{fg=backgr1}

\newlength\subtitwd
\setlength\subtitwd{4cm}% change the width here

\makeatletter
\newcommand\titlegraphicii[1]{\def\inserttitlegraphicii{#1}}
\titlegraphicii{}
\newcommand\superviser[1]{\def\insertsuperviser{Научный руководитель: #1}}
\superviser{}

\setbeamertemplate{title page}
{
  \vbox{}
   {\usebeamercolor[fg]{titlegraphic} \hspace{0.35ex} \inserttitlegraphic\hfill\inserttitlegraphicii \hspace{1ex} \par }\vspace{1.5ex}
  \begin{centering}
    \begin{beamercolorbox}[sep=8pt,center]{institute}
      \usebeamerfont{institute}\insertinstitute
    \end{beamercolorbox}
    \begin{beamercolorbox}[sep=8pt,center]{title}
    
      \usebeamerfont{title}\inserttitle\par%
      \ifx  \insertsubtitle\@empty%
      \else%
        \vskip0.5em%
        {\usebeamerfont{subtitle}\usebeamercolor[fg]{subtitle}\insertsubtitle\par}%
      \fi%     
    \end{beamercolorbox}%
    \vskip1em\par
    \begin{beamercolorbox}[sep=5pt,center]{date}
      \usebeamerfont{date}\insertdate
    \end{beamercolorbox}%\vskip0.5em
    \begin{beamercolorbox}[sep=5pt,center]{author}
      \usebeamerfont{author}\insertauthor
    \end{beamercolorbox}
        \begin{beamercolorbox}[sep=4pt,center]{institute}
      \usebeamerfont{institute}\insertsuperviser
    \end{beamercolorbox}
  \end{centering}
  %\vfill
}
\makeatother

\setbeamercolor{item projected}{bg=ex3}
\setbeamertemplate{enumerate items}[default]

\setbeamercolor{palette primary}{bg=white}
\setbeamercolor{palette primary}{fg=black}
\setbeamercolor{palette secondary}{bg=white}
\setbeamercolor{palette secondary}{fg=black}
\setbeamercolor{palette tertiary}{bg=white}
\setbeamercolor{palette tertiary}{fg=black}

\setbeamercolor{itemize item}{fg=ex3}
\setbeamercolor{itemize subitem}{fg=ex2}
\setbeamercolor{itemize subsubitem}{fg=ex1}

\setbeamercolor{enumerate item}{fg=ex3}
\setbeamercolor{enumerate subitem}{bg=ex3}
\setbeamercolor{enumerate subsubitem}{bg=ex3}


\setbeamertemplate{itemize subitem}{$\Rightarrow$}
\setbeamertemplate{itemize item}{$\blacktriangleright$}



\usepackage{todo}
\newcolumntype{a}{>{\columncolor{red}}c}


\usefonttheme{professionalfonts}

\title[OLG]{Модели пересекающихся поколений (OLG)}


\author[Михаил Гареев]{Михаил Гареев \\ \smallskip \scriptsize ЭО-15-01 \\ \smallskip \scriptsize \href{mailto:mkhlgrv@gmail.com}{\nolinkurl{mkhlgrv@gmail.com} }}

\superviser{к.э.н. Полбин А.В.}

\institute[РАНХиГС]{ \uppercase{
 Российская Академия Народного Хозяйства и  \\ Государственной Службы при Президенте Российской Федерации}}
\date{2020}


\titlegraphic{\includegraphics[scale=0.5]{logo/logo_ranepa.png}}
\titlegraphicii{\includegraphics[scale=0.5]{logo/logo_emit.png}}

\begin{document}

\frame[plain]{\titlepage}	% Титульный слайд


\section{Возникновение и основная идея}

\begin{frame}
 \frametitle{\insertsection}
    \begin{itemize}
        \item Проблема: модели экономического роста первого поколения (модель Солоу) предполагают экзогенную норму сбережения.
        \item Решение: модель, в которой сбережения формируются исходя из выбора экономических агентов:
        \begin{itemize}
        \item модель Рамсея --- Кааса --- Купманса (1963): бесконечно живущий агент выбирает уровень потребления и сбережения, максимизируя полезность;
        \item модель пересекающихся поколений Даймонда (1965): в каждый момент времени существует несколько поколений, молодые работают и сберегают, пожилые тратят свои сбережения. 
        \end{itemize}
    \end{itemize}
\end{frame}

\section{Базовая модель 2-х поколений}
\subsection{Предположения}
\begin{frame}
 \frametitle{\insertsection}
 \framesubtitle{\insertsubsection}
     \begin{itemize}
        \item дискретное время,
        \item индивиды живут два периода,
        \item полезность индивида $U_t = \ln(c_{1,t}) + \beta \ln(c_{2, t+1 })$,
        \begin{itemize}
            \item $c_{1,t}$ --- потребление молодого поколения,
        \item $c_{2,t + 1}$ --- потребление пожилого поколения,
        \item $\beta \in (0, 1)$ --- коэффициент временных предпочтений,
        \end{itemize}
        \item численность населения растёт с постоянным темпом $n$: $L_t = (1+n)^t L_0$,
        \item индивиды работают только в первый период жизни, неэластично предлагая единицу труда и получая зарплату $w_t$,
        
    \end{itemize}

\end{frame}

\begin{frame}
 \frametitle{\insertsection}
 \framesubtitle{\insertsubsection}
     \begin{itemize}
     \item выпуск описывается функцией Кобба-Дугласа: $Y_t =F(K_t, L_t) = {K_t}^\alpha {L_t}^{1-\alpha}$
        \item рынки факторов производства находятся в совершенной конкуренции,
        \item обозначим $k_t = K_t/L_t$, $f(k_t) = F(K_t, 1)$, тогда рентная цена капитала, равная доходности сбережений $R_t = 1-\delta + f^{'}(k_t)$, где $\delta$ --- амортизация,
        \item зарплата $w_t = f(k_t) - R_t k_t$,
        \item единица сбережений $s_t$ без потерь трансформируется в единицу капитала $K_t$.
    \end{itemize}

\end{frame}

\subsection{Задача потребителя}
\begin{frame}
 \frametitle{\insertsection}
 \framesubtitle{\insertsubsection}
 \begin{itemize}
     \item Уровень сбережений $s_t$ индивида из поколения $t$ --- это результат решения оптимизационной задачи:
     \begin{align*}
  \max_{c_{1, t}, {c_{2, t+1}, s_t}} \ln(c_{1,t}) &+ \beta \ln(c_{2, t+1 })\\
  \intertext{при ограничениях:} 
  c_{1,t} + s_t &\leq w_t, \\
c_{2, t} &\leq R_{t+1} s_t.
\end{align*}
\item Условия первого порядка:
\begin{equation*}
    s_t = \frac{\beta}{1+\beta} w_t.
\end{equation*}
 \end{itemize}
 

\end{frame}

\subsection{Динамика капитала}
\begin{frame}
 \frametitle{\insertsection}
 \framesubtitle{\insertsubsection}
 \begin{itemize}
     \item Уровень капитала в момент времени $t+1$:
     \begin{equation*}
  K_{t+1} &= s_t L_t + (1-\delta) K_t.
  \end{equation*}
  \item Перейдём к капиталовооружённости на единицу труда:
  \begin{align*}
       k_{t+1} &= \frac{s_t}{1+n} + (1-\delta) \frac{k_t}{1+n};\\
  k_{t+1} &= \frac{\beta w_t}{(1+\beta)(1+n)} + (1-\delta) \frac{k_t}{1+n};\\
  k_{t+1} &=  \frac{\beta (1-\alpha) k_{t}^\alpha + (1-\delta)  k_t}{(1+\beta)(1+n)}.
  \end{align*}
 
  \item В устойчивом состоянии капиталовооружённость не меняется: $k^{*} = k_{t+1} = k_{t}$:
  \begin{equation*}
  k^{*} &= \left[\frac{\beta (1-\alpha)}{(1+\beta)(1+n) - (1-\delta)}\right]^{\left(\frac{1}{1-\alpha}\right)}
 \end{equation*}
\end{itemize}
 

\end{frame}

\subsection{Динамическая неэффективность}


\section{Модель 3-х поколений}
\subsection{Предпосылки}
\subsection{Решение}
\subsection{Симуляция}
% усложнение модели
% вечная молодость 
% эндогенная рождаемость
% многомерное население + симуляция


\end{document}